\documentclass[10.5pt,a4paper]{article}
\usepackage{fontspec}
\setmainfont{Avenir Next}
\newfontfamily\display{Avenir Next Condensed}[UprightFont=* Demi Bold]
\setmonofont{Menlo}[Scale=0.85]
\usepackage{geometry}
\geometry{margin=2cm,top=1.8cm,bottom=2cm}
\usepackage{xcolor}
\definecolor{ink}{HTML}{1A1A1A}
\definecolor{rule}{HTML}{9A9A9A}\definecolor{RULE}{HTML}{9A9A9A}
\definecolor{red3v}{HTML}{C0392B}
\definecolor{gnd}{HTML}{333333}
\definecolor{sig}{HTML}{1F5FAD}
\usepackage[nobottomtitles*]{titlesec}
\titleformat{\section}{\display\Large\MakeUppercase}{\thesection}{0.8em}{}[\vspace{-0.6ex}{\color{rule}\rule{\linewidth}{0.4pt}}]
\titleformat{\subsection}{\display\large\MakeUppercase}{\thesubsection}{0.8em}{}
\titlespacing*{\section}{0pt}{2.2ex}{1.4ex}
\titlespacing*{\subsection}{0pt}{1.8ex}{0.8ex}
\usepackage{enumitem}
\usepackage{needspace}
\setlist{nosep,leftmargin=1.4em,topsep=0.4ex}
\usepackage{booktabs,tabularx,array}
\renewcommand{\arraystretch}{1.25}
\usepackage{hyperref}
\hypersetup{colorlinks,linkcolor=ink,urlcolor=sig}
\usepackage{tikz}
\usetikzlibrary{positioning,calc}
\usepackage[siunitx]{circuitikz}
\usepackage{fancyhdr}
\pagestyle{fancy}\fancyhf{}
\fancyhead[L]{\small\display SWALLOWCUE PREP PACK}
\fancyhead[R]{\small Oxford Neurotech Hackathon, 7 to 8 Nov 2026}
\fancyfoot[R]{\small\thepage}
\renewcommand{\headrulewidth}{0.3pt}
\setlength{\parindent}{0pt}
\setlength{\parskip}{0.6ex}
\newcommand{\f}[1]{\texttt{#1}}
\newcolumntype{L}{>{\raggedright\arraybackslash}X}

\begin{document}

{\display\Huge SWALLOWCUE}\\[0.3ex]
{\display\Large PREP PACK, OXFORD NEUROTECH HACKATHON}\\[0.5ex]
Sat 7 Nov 09:00 to Sun 8 Nov 18:00. Mathematical Institute, Andrew Wiles Building, Woodstock Road, Oxford OX2 6GG.\\
Written Fri 6 Nov. Place confirmed (Discord joined, name matches Luma). Tracks, hardware and schedule are announced on Saturday morning.

\vspace{1ex}
\section*{The ten things that matter}
\begin{enumerate}[leftmargin=1.8em]
\item \textbf{Write \f{is\_swallow()} tonight}, in \f{warmup/swallow.py} (Python, for EMG) and \f{planb\_imu/src/is\_swallow.cpp} (C++, for the throat sensor). Until then nothing counts as a swallow.
\item \textbf{Leave Marble Arch by 06:15} on the Oxford Tube (student return with the organisers' code). Registration opens 08:45, talks at about 09:45.
\item \textbf{Bring your own kit} (section \ref{sec:kit}). Hardware is a surprise; the throat-sensor version works with no event hardware at all.
\item \textbf{First hour: find out what records EMG.} That picks the signal path (section \ref{sec:first}) in five minutes.
\item \textbf{Pick the track SwallowCue fits} with the decision table (section \ref{sec:track}). Decide by 12:00.
\item \textbf{An ugly loop that works beats a pretty half}: signal to buzz end to end by Saturday 18:00.
\item \textbf{Measure, don't claim.} Score real sessions with \f{tools/eval\_session.py}; test on a person you did not tune on (section \ref{sec:measure}).
\item \textbf{Film the demo working} on Saturday night, as a backup.
\item \textbf{Claim only what you measured}: detects swallows and cues in healthy volunteers. Never ``reduces drooling''.
\item \textbf{Be home Sunday night}: MPHY0010 Scenario 3 starts Monday and is compulsory.
\end{enumerate}

\needspace{14\baselineskip}\section{Tonight, Friday 6 November}
\begin{tabularx}{\linewidth}{@{}p{2.3cm}L@{}}
\toprule
By 20:00 & Write \f{is\_swallow()} in \f{warmup/swallow.py}. Run \f{python3 -m pytest -q test\_swallow.py} in \f{warmup/} until 4 of 4 pass. Start with \f{python3 -c "import swallow; swallow.print\_burst\_stats()"}.\\
By 20:30 & Port the same idea to \f{planb\_imu/src/is\_swallow.cpp}. Different signal, so treat the numbers as a first guess: you tune them on the real sensor tomorrow. Add a gyro limit for head nods.\\
By 21:00 & Run the whole EMG pipeline with no hardware:\newline
\f{uv run live\_emg/swallowcue\_live.py -{}-board fake -{}-speed 10 -{}-log /tmp/f.csv}\newline
then \f{python3 tools/eval\_session.py /tmp/f.csv}. It now uses YOUR rule.\\
By 21:30 & Pack (section \ref{sec:kit}). Charge laptop, phone, power bank. Install the Oxford Tube app and buy the student return with the discount code from \#announcements.\\
By 22:00 & Post in \#find-a-team (section \ref{sec:team}). Text Radin. Alarm for 05:15.\\
\bottomrule
\end{tabularx}

\needspace{14\baselineskip}\section{Getting there, Saturday}
\begin{tabularx}{\linewidth}{@{}p{2.3cm}L@{}}
\toprule
05:30 & Leave home. Archway to Marble Arch: Northern line to Tottenham Court Road, Central line 2 stops west (about 30 min). The Night Tube runs on both lines on Friday night into Saturday.\\
06:15 & Oxford Tube from Marble Arch (Park Lane stop, by the north-east corner of Hyde Park). It runs 24 hours, every few minutes at peak, less often early on Saturday. Allow 2 h: morning journeys are 1 h 40 to 2 h 15.\\
about 08:15 & Get off at Gloucester Green (city centre). Walk north up St Giles and Woodstock Road to the Radcliffe Observatory Quarter, about 15 to 20 min (check the map).\\
08:45 & Registration. Show your Luma ticket (Luma app or the email ``My Ticket''). Your Discord name must match Luma.\\
\bottomrule
\end{tabularx}
Faster but dearer: train from Marylebone or Paddington, about 1 h, then 20 min walk from Oxford station. Home on Sunday: the event ends at 18:00; Oxford Tube back from Gloucester Green, about 2 h.

\needspace{14\baselineskip}\section{Kit list}\label{sec:kit}
\begin{tabularx}{\linewidth}{@{}p{4.2cm}L@{}}
\toprule
\textbf{Signal and cue} & ESP32 DevKit (ELEGOO kit), GY-521 motion sensor (plus the kit's GY-6500 as a spare), 2 coin vibration motors, PN2222 (2 in the kit), 1N4007-type diode, 1 k\(\Omega\) resistors, breadboard, 20 jumpers incl.\ 4 long (30 cm+) soft ones, USB cable that carries data (test it tonight), DRV8833 board as spare motor driver.\\
\textbf{Body} & Snap gel electrodes (TIGA-MED, trim to fit under the chin), alcohol wipes, medical tape (micropore), a cheap razor if you have stubble, scissors.\\
\textbf{Wearable} & An old watch strap or a strip of velcro and tape for the wristband; a small zip bag to hold the board.\\
\textbf{Laptop} & Charger, power bank, USB hub or adapter, mouse. Repos and \f{uv} caches already warm (run each command once tonight while on Wi-Fi).\\
\textbf{Tools} & Multimeter, small screwdriver, spare jumpers, a few M3 bolts if you bring a printed holder.\\
\textbf{You} & Water bottle (they asked), hoodie and something to sleep on if you stay (classrooms only, no bedding), toiletries, earplugs, phone charger, student ID, bank card.\\
\bottomrule
\end{tabularx}

\needspace{14\baselineskip}\section{Timeline and checkpoints}
Shift the times once the schedule is announced; keep the order and the gaps.

\begin{tabularx}{\linewidth}{@{}p{2.6cm}L@{}}
\toprule
Sat 08:45 to 09:45 & Register. Find the hardware table. Ask the questions in section \ref{sec:first}. Meet people from \#find-a-team.\\
09:45 to 10:30 & Opening talks: write down every track name, its brief, any stated metric, and the judges.\\
by 11:30 & Team formed (3 to 5). Roles from section \ref{sec:team}.\\
by 12:00 & \textbf{Track and signal path locked.} One sentence: who it is for, what it senses, what it does.\\
by 14:00 & First real signal on screen: chin EMG or throat motion, with a swallow visible in the envelope.\\
by 18:00 & \textbf{Loop works end to end}: signal, burst, \f{is\_swallow()}, buzz on the wrist. Ugly is fine.\\
18:00 to 22:00 & Tune \f{is\_swallow()} on person A. Record scored sessions (section \ref{sec:measure}). Film the working demo.\\
22:00 to 00:00 & Slides v1, dashboard polish. Then sleep 4 to 5 h. Sleep is a feature.\\
Sun 07:00 to 10:00 & Held-out test on person B (not used for tuning). Fill the numbers slide.\\
10:00 to 12:00 & Second feature only if the core is solid (see the never-add list in section \ref{sec:risks}).\\
Freeze 3 h before judging & No new code. Rehearse the pitch twice with the live demo and once with the video.\\
\bottomrule
\end{tabularx}

\needspace{14\baselineskip}\section{First hour at the venue}\label{sec:first}
\subsection{Questions for the organisers}
\begin{enumerate}
\item What hardware can we use, how many units, and can we take one to our table for the whole weekend?
\item Does anything record EMG (OpenBCI Cyton, Ganglion, a muscle sensor)? If it is a Cyton with a Daisy board, can we remove the Daisy (it halves the rate to 125 Hz)?
\item Are electrodes on the face and neck allowed on volunteers? Is there a consent form?
\item Can we use our own hardware and code written before the event (declared)?
\item When and how is judging done: track rooms, demo fair, or main stage? How long is the pitch?
\item What do judges score? (If unstated, assume innovation, technical depth, impact, presentation, 25\% each.)
\end{enumerate}

\subsection{Pick the signal path from what they have}
\begin{tabularx}{\linewidth}{@{}p{4.4cm}p{3.2cm}L@{}}
\toprule
They provide & Path & Run\\
\midrule
OpenBCI Cyton (250 Hz) & A: chin EMG & \f{uv run live\_emg/swallowcue\_live.py -{}-board cyton -{}-port <dongle> -{}-channel 1 -{}-cue-port <esp32> -{}-log s1.csv}\\
Cyton + Daisy & A, Daisy off & Remove the Daisy (250 Hz). If you cannot: \f{-{}-board cyton\_daisy}; it runs at 125 Hz with a 10 to 56 Hz band, weaker.\\
OpenBCI Ganglion (200 Hz) & A & \f{-{}-board ganglion -{}-port ...}; channel 1 is a natively bipolar pair.\\
Analog muscle sensor (MyoWare-style, ENV output) & C: envelope on ESP32 & ENV to GPIO 34, sensor powered from 3V3. Flash \f{planb\_imu}, type \f{e} + Enter. Same detector, same plot.\\
EEG only (Muse, Crown, Unicorn) & B: throat motion & Your GY-521 kit. Nothing from them needed.\\
Nothing usable & B & Same.\\
\bottomrule
\end{tabularx}
In every path the ESP32 is the wristband. With no motion sensor attached it starts in wristband-only mode and buzzes when the laptop sends \f{b}.

\needspace{14\baselineskip}\section{Which track, and how to frame it}\label{sec:track}
The site promises four tracks ``from peripheral nervous-system interfaces to sub-cortical implants''. Match on the brief, not the name.

\begin{tabularx}{\linewidth}{@{}p{3.7cm}p{2.3cm}L@{}}
\toprule
Track type & Fit & Framing\\
\midrule
Peripheral interfaces, wearables, EMG & Best & A peripheral sensor closes a loop on a motor symptom. Lead with the chin EMG and the live buzz.\\
Closed-loop neuromodulation or neurorehab & Strong & Sense, decide, deliver a sensory cue, only when needed. This is the organiser's own field (closed-loop, Denison lab), so name the loop and its safety rules: never cue more often than every T, stop when swallowing resumes.\\
AI, data, decoding & Good with an add-on & Rule vs learned classifier on burst features (duration, peak, gaps, gyro), tested on a person held out of training. Report both, and say which you would trust and why.\\
Non-invasive BCI (EEG) & Weak & Do not force it. If teams must stay in one track, enter the peripheral one. If EEG is mandatory, switch to the runner-up in PREP.md 4b.\\
Implants, DBS, sub-cortical & Weak & Concept tracks are won on framing. Only enter if no better fit: SwallowCue as the sensing layer a DBS clinic could use to log swallow rate.\\
Sponsor track with a metric & Depends & If the brief states a number, build the scoring script in hour 1 and optimise that number.\\
\bottomrule
\end{tabularx}

\needspace{14\baselineskip}\section{What you are building}
\subsection{The loop}
Signal (chin EMG, throat motion, or a muscle-sensor envelope) \(\rightarrow\) remove slow drift and mains hum \(\rightarrow\) rectify or take magnitude \(\rightarrow\) 5 Hz envelope \(\rightarrow\) burst when the envelope is above 3\(\times\) the quiet level \(\rightarrow\) \f{is\_swallow(burst)} \(\rightarrow\) cue timer: buzz \(T\) s after the last swallow, then every \(T\) s until the next one.

Each burst carries duration, peak, gap before, gap after (capped at 2 s, because the device cannot wait forever), and, for the throat sensor, peak head rotation.

\subsection{What is done and tested}
\begin{tabularx}{\linewidth}{@{}p{5.4cm}L@{}}
\toprule
\f{live\_emg/swallowcue\_live.py} & Live EMG pipeline. Tested on fake sessions (with mains hum) and on BrainFlow's synthetic board. Not yet on a real Cyton.\\
\f{planb\_imu/} (ESP32 firmware) & Throat motion, wristband-only and EMG-envelope modes. Builds; detector passes 6 of 6 tests on the laptop. Not yet flashed or worn.\\
\f{tools/live\_plot.py} & Demo screen for every path (serial port or a log file). Tested with fake input.\\
\f{tools/eval\_session.py} & Scores sessions for the slide. Tested on fake sessions.\\
\f{is\_swallow()}, both versions & \textbf{Yours to write.}\\
\bottomrule
\end{tabularx}

\subsection{Command cheat sheet}
\begin{tabularx}{\linewidth}{@{}p{4.3cm}L@{}}
\toprule
Flash the ESP32 & Ask me: \f{pio\_build}, \f{pio\_devices}, \f{pio\_upload} on \f{planb\_imu}.\\
ESP32 serial commands & Type in the live-view terminal, then Enter: \f{r} recalibrate (sit still 3 s), \f{t15} cue period 15 s, \f{b} test buzz, \f{e} EMG-envelope mode on GPIO 34, \f{i} back to motion sensor. A blank Enter marks a swallow.\\
Throat-sensor demo screen & \f{uv run tools/live\_plot.py <esp32 port> -{}-log s1.csv} and press Enter after each swallow.\\
EMG live & \f{uv run live\_emg/swallowcue\_live.py -{}-board cyton -{}-port <p> -{}-cue-port <esp32> -{}-log e1.csv -{}-T 15}\\
EMG demo screen & In a second terminal: \f{uv run tools/live\_plot.py e1.csv}\\
Score & \f{python3 tools/eval\_session.py e1.csv e2.csv -{}-period 15}\\
Find ports & \f{ls /dev/cu.*} (the OpenBCI dongle is an FTDI \f{usbserial}).\\
\bottomrule
\end{tabularx}

\needspace{14\baselineskip}\section{Wiring (paths B and C)}
\begin{center}
\begin{tikzpicture}[x=1cm,y=1cm,font=\small,
  blk/.style={draw=ink,line width=0.6pt,fill=white},
  pn/.style={font=\footnotesize\ttfamily},
  w/.style={line width=0.8pt,sig},
  pw/.style={line width=0.8pt,red3v},
  gw/.style={line width=0.8pt,gnd}]
% ESP32
\draw[blk] (0,0) rectangle (3.2,8);
\node[font=\display] at (1.6,8.35) {ESP32 DEVKIT};
\foreach \y/\n in {7.4/3V3,6.6/GND,5.8/GPIO21 SDA,5.0/GPIO22 SCL,4.0/GPIO34 ADC,1.6/GPIO25}{
  \node[pn,anchor=east] at (3.1,\y) {\n};
  \fill[ink] (3.2,\y) circle (1.4pt);
}
% GY-521
\draw[blk] (8.4,4.6) rectangle (10.8,8);
\node[font=\display] at (9.6,8.35) {GY-521};
\foreach \y/\n in {7.4/VCC,6.6/GND,5.8/SDA,5.0/SCL}{
  \node[pn,anchor=west] at (8.5,\y) {\n};
  \fill[ink] (8.4,\y) circle (1.4pt);
}
\node[font=\footnotesize,align=left,anchor=west] at (11.0,6.4) {taped to\\the throat,\\long soft\\wires};
\draw[pw] (3.2,7.4) -- (8.4,7.4);
\draw[gw] (3.2,6.6) -- (8.4,6.6);
\draw[w] (3.2,5.8) -- (8.4,5.8);
\draw[w] (3.2,5.0) -- (8.4,5.0);
% analog envelope input (path C)
\draw[w] (3.2,4.0) -- (4.2,4.0);
\node[font=\footnotesize,anchor=west] at (4.25,4.0) {ENV in, path C only};
% motor stage
\draw (3.2,1.6) to[R, l={1 kΩ}] (5.4,1.6) -- (6.0,1.6);
\draw (6.0,1.6) node[npn,anchor=B](Q){};
\node[font=\footnotesize,anchor=west] at ($(Q.B)+(0.85,-0.05)$) {PN2222};
\node[font=\footnotesize,anchor=east] at ($(Q.C)+(-0.08,0.12)$) {C};
\node[font=\footnotesize,anchor=east] at ($(Q.E)+(-0.08,-0.12)$) {E};
\draw (Q.E) -- ++(0,-0.3) node[ground]{};
\draw (Q.C) |- (7.2,2.7) coordinate (m1);
\draw (m1) -- (9.4,2.7) coordinate (m2);
\node[draw=ink,circle,fill=white,minimum size=0.6cm,font=\footnotesize] at (8.3,2.7) {M};
\node[font=\footnotesize,anchor=north] at (8.3,2.35) {coin motor};
\draw (m1) -- (7.2,3.3) to[D] (9.4,3.3) -- (m2);
\node[font=\footnotesize,anchor=south west] at (8.55,3.45) {stripe $\rightarrow$};
\draw[pw] (9.4,3.0) -- (10.3,3.0) node[right,font=\footnotesize\ttfamily,text=red3v] {to 3V3};
\end{tikzpicture}
\end{center}
\textbf{Check the PN2222 leg order} on the kit card before powering: some TO-92 parts are E-B-C, others C-B-E. The motor draws about 63 mA from 3V3. Put tape between the GY-521's back and the skin: its solder joints are bare.

\needspace{14\baselineskip}\section{Putting it on a person}
\subsection{Chin EMG (path A)}
\begin{enumerate}
\item Shave and wipe the skin under the chin with alcohol; let it dry.
\item Two gel electrodes under the chin, either side of the midline, centres about 2 cm apart, over the submental muscles (the floor of the mouth). Trim the 43 mm foam if they overlap.
\item Reference and bias on bone: behind the ear (mastoid) or on the collarbone.
\item Cyton, simplest: first chin electrode to \f{N1P}, second to \f{SRB2}, bone electrode to \f{BIAS}. Or use \f{N1P}/\f{N1N} with \f{-{}-bipolar} (it sends the channel setting \f{x1060100X}; check that string against the OpenBCI docs on the day before relying on it).
\item Check: clench the jaw, swallow, talk. You should see three different bursts in the plot. If the trace is a fat 50 Hz band, an electrode is loose.
\item Tape the leads to the collar so a head turn does not pull them.
\end{enumerate}

\subsection{Throat motion (path B)}
Flat on the front of the neck, just below the Adam's apple, chip side out, held with micropore tape or a loose elastic band. Wires run down to the collar. Too tight changes how swallowing feels; too loose and talking rattles it. Sit still for 3 s after power-up (calibration).

\needspace{14\baselineskip}\section{Measuring: the number on the slide}\label{sec:measure}
PI judges reward a measured number on a real signal. Plan for it from Saturday afternoon.

\subsection{Protocol, about 6 min per person}
\begin{tabularx}{\linewidth}{@{}p{3.2cm}p{2cm}L@{}}
\toprule
Block & Length & What happens\\
\midrule
Quiet sips & 1 min & 5 sips of water, about 10 s apart. Press Enter right after each swallow.\\
Reading aloud & 1 min & Read a paragraph. No swallowing if you can help it; mark any you do.\\
Chewing gum & 1 min & Chew normally; mark swallows.\\
Conversation & 1 min & Talk with a teammate; mark swallows.\\
Head movement & 1 min & Nod and turn, 10 times each (throat sensor: tests the gyro rule).\\
Hold off & 1 min & Do not swallow until the buzz, then swallow (tests the cue).\\
\bottomrule
\end{tabularx}
Save each block as its own file (\f{-{}-log A\_read.csv}, \f{A\_chew.csv}, ...), so the score splits by activity.

\subsection{Honest testing}
\begin{itemize}
\item Tune \f{is\_swallow()} on person A only. Freeze it. Then record person B and report B's numbers as the result. Say so on the slide.
\item Report: swallows caught (\%), false counts per minute per activity, and buzzes vs a fixed schedule with the same period. \f{eval\_session.py} prints all of these.
\item State N (people, swallows, minutes). Small N is fine if it is stated.
\item What to expect: an offline under-chin EMG study in Parkinson's reported 83\% (MDPI Sensors 2025). Healthy volunteers in a weekend will be different; do not compare directly.
\end{itemize}

\needspace{14\baselineskip}\section{Pitch}
\subsection{Structure, 3 minutes}
\begin{tabularx}{\linewidth}{@{}p{2cm}p{2.6cm}L@{}}
\toprule
Time & Slide & Point to land\\
\midrule
0:00 to 0:30 & The person & Drooling in Parkinson's is common (often over half of people in studies) and isolating. Saliva is normal; swallowing is too rare. Treatments are anticholinergic drugs (side effects) or botulinum toxin injections.\\
0:30 to 1:00 & What exists & Cueing works: a once-a-minute wristband reduced drooling (McNaney 2019, 28 people). CueBand is in a 300-person trial (Parkinson's UK). Both buzz on a fixed timer, whether or not you just swallowed.\\
1:00 to 1:50 & Live demo & Talk: nothing. Sip: count +1. Chew or nod: nothing. Hold off: buzz. Swallow: reset.\\
1:50 to 2:30 & The number & Held-out person: \% caught, false counts per minute by activity, buzzes vs fixed timer. N stated.\\
2:30 to 3:00 & What next & Clinician-set \(T\), a swallow-rate log for the neurologist, a pilot with people with Parkinson's. One line on cost (parts under GBP 30).\\
\bottomrule
\end{tabularx}
Write the words yourself; this is the skeleton. Practise the demo order until the transitions are automatic.

\subsection{Positioning}
Not ``the first''. Northwestern's Aria patch (trial NCT04664634, about 60 people, no results posted) already senses swallows and cues. Swallowing accelerometry is an established research method. Your line: a cheap, open, smart version of the cue that people already wear, which buzzes only when swallowing stops, measured live.

\needspace{14\baselineskip}\section{Questions judges will ask}
\begin{tabularx}{\linewidth}{@{}p{5.2cm}L@{}}
\toprule
Why not just CueBand? & A fixed timer buzzes when you have just swallowed, so people learn to ignore it; ours buzzes only when needed. Show buzzes vs the fixed schedule from your session.\\
How do you tell a swallow from talking or chewing? & Duration plus the quiet around it (speech is many short bursts close together, chewing repeats every 0.7 s), plus head rotation for the motion sensor. Show the burst table.\\
Which error is worse, a missed swallow or a false one? & A missed swallow costs one extra buzz (annoying). A false swallow silences a needed cue (drooling). So the rule leans towards missing rather than inventing swallows. Your call: be ready to defend the one you made.\\
Does it work in Parkinson's? & Not tested. Healthy volunteers only. Parkinson's swallowing is slower and weaker, so thresholds would be set per person in clinic.\\
Why EMG and not a microphone? & EMG (or motion) does not hear other people talking and is not affected by room noise; Aria went acoustic, which is the comparison.\\
Beards, sweat, electrodes falling off? & Shaved skin today; long-term it would be a dry-electrode chin patch or the motion sensor. The motion path is the more wearable one.\\
Battery and size? & ESP32 is a dev board here; a production version is a coin-cell patch plus a band, BLE between them.\\
Is it a medical device? & Probably yes in the UK if it claims to treat drooling (MHRA); a low-risk class is likely since it only vibrates, but that needs checking. Today it is a research prototype.\\
What would you do with GBP 5k? & Build 5 units, run a 2-week pilot with a Parkinson's clinic and compare swallow rate and drooling score against the fixed-timer cue.\\
What did the team build this weekend vs before? & Say exactly: the warm-up code and firmware were written before (declared); the tuning, measurements, wearable and demo happened here.\\
\bottomrule
\end{tabularx}

\needspace{14\baselineskip}\section{Risks and fixes}\label{sec:risks}
\begin{tabularx}{\linewidth}{@{}p{5.2cm}L@{}}
\toprule
No EMG hardware & Path B, already built.\\
Cyton at 125 Hz (Daisy) & Remove the Daisy, or run \f{cyton\_daisy}; envelope still works.\\
50 Hz everywhere & Notch is on by default; check electrode contact; laptop on battery.\\
Talking counts as swallows & Raise the quiet-around requirement; add duration window; measure per activity.\\
Nods count (motion sensor) & Gyro limit in \f{is\_swallow()}.\\
ESP32 not seen by the laptop & Data cable; hold BOOT while flashing; try the GY-6500 board if the GY-521 is dead.\\
Live demo fails on stage & Play the filmed backup and say it is a recording.\\
Overrun & Feature freeze 3 h before judging. The never-add list: phone app, cloud, ML model you cannot explain.\\
\bottomrule
\end{tabularx}

\needspace{14\baselineskip}\section{Team}\label{sec:team}
\begin{itemize}
\item \textbf{You}: signal processing, \f{is\_swallow()}, measurement.
\item \textbf{Hardware}: wristband, wiring, wearable comfort.
\item \textbf{Screen and data}: dashboard, logging, the numbers slide.
\item \textbf{Story}: the patient, the evidence, the pitch. Ideally a medic, neuroscience or speech-and-language student.
\end{itemize}
In \#find-a-team, post three lines in your own words: what you will build (a wearable that cues swallowing in Parkinson's when it stops), what you bring (signal processing, ESP32 firmware, already a working pipeline), who you need (someone clinical, someone for slides and story). Text Radin before you post.

Credit afterwards is team credit. Never present it as solo.

\needspace{14\baselineskip}\section{Safety and ethics}
\begin{itemize}
\item Volunteers are teammates who agree, and can stop any time. No data leaves the laptop; delete recordings after the event unless everyone agrees to keep them.
\item The only output to the body is a vibration motor. No electrical stimulation.
\item Laptop on battery while electrodes are on a person. OpenBCI boards run on their own battery; the ESP32 never touches skin electrically (tape under the GY-521).
\item Stop if skin reddens under an electrode.
\end{itemize}

\needspace{14\baselineskip}\section{After the event}
\begin{itemize}
\item Within 2 days: a one-line thank-you to each judge, with the repo or video link.
\item Push the code with the team's names; keep \f{is\_swallow()} credited to whoever wrote it.
\item Post the 30 s video. It is the first filmed, worn build for your portfolio: claim only what it shows.
\end{itemize}

\needspace{14\baselineskip}\section{Sources}
\small
\begin{itemize}
\item Event: \url{https://oxneurotech.com/hackathon}, \url{https://luma.com/ceww7ooe}, the event Discord (\#start-here, \#announcements, \#questions, read 6 Nov).
\item Drooling in Parkinson's, prevalence and mechanism: \url{https://pmc.ncbi.nlm.nih.gov/articles/PMC7669801}; Cochrane protocol on botulinum toxin: \url{https://www.cochranelibrary.com/cdsr/doi/10.1002/14651858.CD012408}.
\item Cueing: McNaney et al.\ 2019 (PD-CUE); CueBand trial ISRCTN59284050; Aria patch NCT04664634. Details and links in \f{RESEARCH\_2026-11-04.md}.
\item OpenBCI sample rates: BrainFlow board descriptions (Cyton 250 Hz, Cyton + Daisy 125 Hz, Ganglion 200 Hz), checked 6 Nov.
\item Oxford Tube: \url{https://bustimes.org/services/tube-oxford-london}.
\end{itemize}

\end{document}
